% =============================================================================
%  1-Introduccion.tex — Contexto y problema que aborda el taller
% =============================================================================

\section{Introducción}

El análisis de la movilidad en la Zona Metropolitana del Valle de México (ZMVM)
suele enfrentarse a una barrera crítica: la fragmentación institucional y
geográfica de los datos espaciales.  Mientras que los flujos de transporte masivo
ignoran las fronteras políticas, las bases de datos gubernamentales y los archivos
GTFS oficiales suelen interrumpirse drásticamente entre la Ciudad de México y el
Estado de México, omitiendo además sistemas periféricos clave como el Mexibús o
rutas complejas del Trolebús \citep{semovi2020diagnostico}.  Para geógrafos,
urbanistas y analistas del territorio, limpiar y corregir estas discontinuidades
geométricas exige habilidades avanzadas de programación, lo que limita la
democratización efectiva de la información urbana.

Este taller práctico de 2~horas propone derribar esa barrera técnica enseñando a
los asistentes a consumir e integrar datos geoespaciales estructurados en tiempo
real utilizando \textbf{Apimetro} \citep{garibaldi2025apimetro}, una API RESTful
de software libre que unifica las redes multimodales de la metrópoli empleando
OpenStreetMap (OSM) como fuente geométrica de respaldo y validación cruzada
\citep{haklay2010good}.  Los participantes aprenderán a explotar esta
infraestructura directamente desde \textbf{QGIS}, transformando consultas
dinámicas a la API en insumos analíticos sin escribir una sola línea de código.
